Jérusalem, céleste cité,\\
Bienheureuse vision de la paix,\\
Bâtie de pierres vivantes\\
Vous vous élevez jusqu'aux astres,\\
Entourée de milliers d'Anges\\
Qui vous font un cortège d'épousée.\\
\\
Dotée par le Père de sa gloire,\\
La grâce de l'Epoux est sur vous répandue ;\\
Reine de toute beauté,\\
Que le Christ Roi s'est unie ;\\
Combien heureux est votre sort,\\
Resplendissante cité des cieux !\\
\\
Faites de perles brillantes,\\
Vos portes demeureront ouvertes pour tous ;\\
Car c'est vers elles que la vertu\\
Conduit le mortel qui la prend pour guide,\\
Quiconque pressé de l'amour du Christ\\
Supporte ici-bas des tourments.\\
\\
Il faut que toute pierre, pour entrer dans vos murs\\
Se livre à l'ouvrier qui la polit\\
Sous les coups répétés du marteau,\\
Du ciseau salutaire ;\\
Il faut qu'elle s'appareille et se laisse fixer\\
Pour y trouver place honorable.\\
\\
Que l’honneur dû au Père\\
Soit en tout lieu donné au Très-Haut\\
Ainsi qu’au Fils unique du Père\\
Et à l’illustre Paraclet,\\
Auquel louange, puissance et gloire\\
Soient dans les siècles éternels.